\chapter{Wstęp}\label{chap:wprowadzenie}

\chapter{Stan wiedzy}\label{chap:stan_wiedzy}
Najwcześniejsze systemy telewizji przemysłowej \textit{(ang. Closed Circuit Television -- CCTV)}
powstały w pierwszej połowie XX wieku i umożliwiały jedynie bieżącą obserwację zdarzeń, ze względu
na brack technologii pozwalających na rejestrowanie i przechowywanie sygnału wideo. Rozwój nośników
szpulowych umożliwił rejestrowanie materiału z monitoringu. Systemy te wymagały były jednak
uciążliwe w obsłudze - wymagały ręcznej wymiany taśm magnetycznych oraz ręcznego nawijania taśmy na
szpulę odbiorczą. Wraz z postępem technologicznym metody zapisu ulegały systematycznej poprawie, co
przyczyniło się do popularyzacji systemów bezpieczeństwa. Już w latach 90. XX opracowano
technologię multipleksowania cyfrowego, umożliwiającą jednoczesne nagrywanie z kilku kamer, a także
nagrywanie w trybie poklatkowym i zapis wyłącznie przy wykryciu ruchu \cite{kreugle}.

Istotny przełom w rozwoju systemów wideonadzoru przyniosła wizja komputerowa, rozwijana od lat 70.
XX wieku, której nadrzędnym celem stało się pełne zautomatyzowane zrozumienie analizowanej sceny.
Prawdziwą rewolucję w tej dziedzinie umożliwiło jednak dopiero znaczne zwiększenie mocy
obliczeniowej, rozwój algorytmów uczenia maszynowego oraz powstanie wielkich, oznaczonych zbiorów
danych (m.in. ImageNet, Microsoft COCO czy LVIS). Bazy te dostarczyły niezbędnych informacji do
trenowania zaawansowanych modeli oraz pozwoliły na obiektywne mierzenie postępów w dziedzinie
rozpoznawania i segmentacji obrazu \cite{szeliski2022}.

%%%%%%%%%%%%%%%%%%%%%%%%%%%%%%%%%%%%%%%%%%%%%%%%%%%%%%%%%%%
%%%%%%%%%%%%%%%%%%%%%%%%%%%%%%%%%%%%%%%%%%%%%%%%%%%%%%%%%%%

\section{Analiza obrazu i podstawowe definicje}

\subsection{Komputerowa reprezentacja obrazu}
Komputer, analizując obraz, sprowadza go do postaci macierzy pikseli, którym przypisane są wartości
liczbowe zależne od koloru. Kształt tej macierzy zależy od przestrzeni barw, w jakiej
reprezentowany jest dany obraz. W przypadku skali szarości obraz jest przedstawiany w postaci
macierzy dwuwymiarowej, której elementy przyjmują wartości od 0 do 255, gdzie 0 oznacza czerń, a
255 biel \cite{gonzalezwoods}.

Obrazy przedstawione w pełnym spektrum kolorów zapisywane są najczęściej w przestrzeni RGB
\textit{(ang. Red, Green, Blue)} lub HSV \textit{(ang. Hue, Saturation, Value)}. W takim przypadku
obraz opisywany jest za pomocą trzech macierzy, a każdemu pikselowi przypisywane są trzy wartości z
zakresu od 0 do 255. W przestrzeni RGB, która jest wykorzystywana częściej, wartości te odpowiadają
odpowiednio natężeniu koloru czerwonego, zielonego i niebieskiego \cite{gonzalezwoods}. Takie
przedstawienie pozwala na odwzorowanie pełnego spektrum barw w formie dogodnej do przetwarzania
komputerowego.

Liczba pikseli tworzących obraz jest równoznaczna z jego rozdzielczością\cite{gonzalezwoods}, a
więc bezpośrednio wpływa na poziom szczegółowości widocznych obiektów i możliwości ich poprawnej
detekcji. Nagrania wideo składają się z pojedynczych obrazów nazywanych w tym kontekście klatką lub
ramką \textit{(ang. frame)}. Ciąg kolejnych klatek wyświetlanych w krótkich odstępach czasu tworzy
wrażenie ruchu i stanowi strumień wideo. Liczba klatek na sekundę, wpływa na płynność nagrania oraz
możliwość dokładnej analizy ruchu obiektów w czasie.

\subsection{Model kamery i rzutowanie na plan terenu}
Dla poprawnej analizy nagrań wideo należy wziąć pod uwagę sposób ich powstawania w kamerze. Aby
możliwe było poprawne wyznaczanie położenia obiektów w scenie, stosuje się proces kalibracji
kamery, który pozwala określić zależność między współrzędnymi punktów w przestrzeni rzeczywistej a
ich położeniem w obrazie. Podstawą takiego opisu jest model otworkowy kamery \textit{(ang. pinhole
  camera model)}, w którym obraz powstaje jako rzut perspektywiczny punktów trójwymiarowej sceny na
dwuwymiarową płaszczyznę obrazu\cite{hartleyzisserman}.

Pełny model kamery można zapisać w postaci iloczynu macierzy wewnętrznej oraz macierzy zewnętrznej.
Macierz wewnętrzna opisuje właściwości samej kamery i sposób odwzorowania punktów w układzie kamery
na współrzędne obrazu. Zawiera ona między innymi ogniskowe w osiach poziomej i pionowej, położenie
punktu głównego oraz ewentualny parametr skośności osi. Parametry te określają więc geometryczne
cechy obrazowania i wpływają na skalę oraz położenie obrazu w układzie pikseli\cite{szeliski2022}.

Grupa parametrów zewnętrznych, opisuje położenie i orientację kamery względem układu współrzędnych
sceny. Obejmują one macierz obrotu oraz wektor translacji. Parametry te opisują transformację
punktów ze współrzędnych świata do układu współrzędnych kamery. Dzięki temu możliwe jest
określenie, z którego miejsca kamera obserwuje scenę oraz w jakim kierunku jest
skierowana\cite{szeliski2022}.

Złożenie macierzy wewnętrznej i zewnętrznej daje pełną macierz projekcji kamery \(\mathbf{P} =
\mathbf{K}[\mathbf{R} \mid \mathbf{t}]\), która opisuje kompletną transformację punktu z
trójwymiarowej przestrzeni sceny do dwuwymiarowego układu pikselowego obrazu
\cite{hartleyzisserman}.

W praktyce rzeczywisty obraz może odbiegać od idealnego modelu z powodu dystorsji soczewki, która
powoduje zniekształcenia, są one szczególnie widoczne na obrzeżach obrazu i wynikają z właściwości
układu optycznego. Z tego względu w procesie kalibracji często wyznacza się również współczynniki
pozwalające na korekcję tych zniekształceń przed dalszą analizą obrazu\cite{szeliski2022}.

Do analizy wielu obrazów tej samej powierzchni lub terenu służy analiza homografii. W dziedzinie
widzenia komputerowego dowolne dwa obrazy tej samej płaskiej powierzchni w przestrzeni są powiązane
przekształceniem rzutowym, zwanym właśnie homografią. Opisuje się ją za pomocą macierzy
\(\mathbf{H}\) o wymiarach \(3 \times 3\), która odwzorowuje punkty z jednej płaszczyzny na drugą z
uwzględnieniem zniekształceń perspektywy. Macierz ta działa na punkty zapisane we współrzędnych
jednorodnych i przyjmuje postać:
\[
  \lambda
  \begin{bmatrix}
    x' \\
    y' \\
    1
  \end{bmatrix}
  =
  \mathbf{H}
  \begin{bmatrix}
    x \\
    y \\
    1
  \end{bmatrix},
\]
gdzie \((x, y)\) oznacza współrzędne punktu w pierwszym układzie, \((x', y')\) współrzędne punktu
po przekształceniu, natomiast \(\lambda\) jest współczynnikiem skali wynikającym z użycia
współrzędnych jednorodnych. Taki zapis pozwala ująć w jednej macierzy przesunięcie, obrót, zmianę
skali oraz zniekształcenie perspektywiczne \cite{zhang2000, szeliski2022}.

W kontekście systemów wielokamerowych rzutowanie na plan terenu umożliwia przedstawienie
trajektorii obiektów w jednej, wspólnej przestrzeni geometrycznej, niezależnie od położenia i
orientacji poszczególnych kamer.

Należy jednak zaznaczyć, że poprawność takiego rzutowania zależy od jakości kalibracji kamery oraz
od tego, w jakim stopniu obserwowana scena spełnia założenie o płaskości terenu. W przypadku
nierównego podłoża lub obiektów o znacznej wysokości konieczne jest zastosowanie bardziej
rozbudowanych modeli geometrycznych.

\subsection{Detekcja ruchu i odejmowanie tła}

Technika odejmowania tła \textit{(ang. background subtraction)} umożliwia rozpoznanie ruchomych
obiektów na nagraniach wideo dla statycznych scen poprzez oddzielenie pierwszego planu od tła.
Przyjmując założenie, że tło pozostaje wzglęgnie niezmienne (mogą pojawiać się zmiany oświetlenia,
zabrudzenia itp.), możliwe jest opisanie tła za pomocą modelu statystyczngo. Stworzony model tła
odejmowany jest od obrazu, pozostawiając jedynie anomalie, czyli obiekty znajdujące się w ruchu,
pierwszy plan zwracany jest zazwyczaj w postaci maski.\cite{zivkovic2004}

\begin{figure}[H]
  \centering
  \includegraphics[width=.8\textwidth]{background-substraction.jpg}
  %\vspace*{-10mm}
  \caption{Wizualizacja procesu odejmowania tła. Źródło: \cite{background_subtraction}}
  \label{fig:mog2}
\end{figure}

W kontekście projektowania zaawansowanych systemów śledzenia wielu obiektów i reidentyfikacji,
algorytmy odejmowania tła pozwalają na ograniczenie obszaru poszukiwań, system może przetwarzać
wyłącznie wykadrowane obszary ruchu (ang. Regions of Interest – ROI).

\subsubsection{MOG2 (Mixture of Gaussians 2)}
Algorytm MOG2 stanowi bezpośrednie rozwinięcie metody mieszanin rozkładów Gaussa zaproponowanej
przez Stauffera i Grimsona. Model GMM \textit{(ang.~Gaussian Mixture Model)} zakłada stałą liczbę
komponentów gaussowskich dla każdego piksela w scenie, co może prowadzić do namdmiernej złożoności
obliczeniowej dla jednorodnych i statycznych obszarów obrazu. MOG2, opracowany przez Zivkovica,
wprowadza adaptacyjny dobór liczby rozkładów dla każdego piksela z osobna.\cite{zivkovic2004}

W~algorytmie MOG2 prawdopodobieństwo zaobserwowania określonej wartości piksela $\vec{x}$ w~chwili
$t$ w~wielowymiarowej przestrzeni barw (np. RGB) jest opisywane przez ważoną sumę $M$ rozkładów
normalnych~\cite{stauffer1999, zivkovic2004}:

\begin{equation}
  \hat{p}(\vec{x} \mid X_T, \text{BG}+\text{FG}) = \sum_{m=1}^{M} \hat{\pi}_m \mathcal{N}(\vec{x}; \hat{\vec{\mu}}_m, \hat{\sigma}_m^2 \mathbf{I})
  \label{eq:mog2_prob}
\end{equation}

gdzie $\hat{\pi}_m$ oznacza oszacowaną wagę $m$-tego komponentu, $\hat{\vec{\mu}}_m$ jest jego
wektorem średnich, natomiast $\hat{\sigma}_m^2 \mathbf{I}$ reprezentuje macierz
kowariancji~\cite{stauffer1999, zivkovic2004}. Dla uproszczenia obliczeń zakłada się, że
poszczególne kanały kolorów są od siebie niezależne i~posiadają tę samą wariancję, co sprowadza
macierz kowariancji do iloczynu wariancji $\hat{\sigma}_m^2$ oraz macierzy jednostkowej
$\mathbf{I}$~\cite{stauffer1999, zivkovic2004}. Wagi mieszaniny są nieujemne i~podlegają
normalizacji, spełniając warunek $\sum_{m=1}^{M} \hat{\pi}_m = 1$~\cite{stauffer1999,
  zivkovic2004}.

Dla każdej nowej klatki obrazu i~każdej nowo zaobserwowanej wartości piksela $\vec{x}(t)$
realizowana jest rekurencyjna aktualizacja parametrów modelu~\cite{stauffer1999, zivkovic2004}.
Najpierw następuje porządkowanie komponentów według stosunku wagi do odchylenia standardowego
($\hat{\pi}_m / \hat{\sigma}_m$), co faworyzuje rozkłady o~wysokiej wadze i~niskiej wariancji jako
najbardziej prawdopodobne definicje tła~\cite{stauffer1999}. Dopasowanie próbki do istniejących
komponentów określa się na podstawie odległości Mahalanobisa~\cite{stauffer1999, zivkovic2004}.
Jeśli kwadrat tej odległości jest mniejszy od przyjętego progu (odpowiadającego zazwyczaj trzem
odchyleniom standardowym), komponent uznaje się za dopasowany~\cite{stauffer1999}. Zmienna
przynależności $o_m^{(t)}$ przyjmuje wówczas wartość 1 dla najlepiej dopasowanego rozkładu i~0 dla
pozostałych~\cite{stauffer1999}. Aktualizacja parametrów zachodzi według następujących równań
rekurencyjnych~\cite{stauffer1999, zivkovic2004}:

\begin{equation}
  \hat{\pi}_m \leftarrow \hat{\pi}_m + \alpha (o_m^{(t)} - \hat{\pi}_m)
  \label{eq:mog2_weight_update}
\end{equation}
\begin{equation}
  \hat{\vec{\mu}}_m \leftarrow \hat{\vec{\mu}}_m + o_m^{(t)} \left( \frac{\alpha}{\hat{\pi}_m} \right) \vec{\delta}_m
  \label{eq:mog2_mean_update}
\end{equation}

\begin{equation}
  \hat{\sigma}^2_m \leftarrow \hat{\sigma}^2_m + o_m^{(t)} \left( \frac{\alpha}{\hat{\pi}_m} \right) \left( \vec{\delta}_m^T \vec{\delta}_m - \hat{\sigma}^2_m \right)
  \label{eq:mog2_var_update}
\end{equation}

gdzie $\vec{\delta}_m = \vec{x}(t) - \hat{\vec{\mu}}_m$, natomiast $\alpha$ określa współczynnik
decydujący o tempie zmian parametrów rozkładu, $\alpha = 1/T$, gdzie $T$ jest rozważanym interwałem
czasowym (liczbą klatek)~\cite{stauffer1999, zivkovic2006}. Jeśli dla danego piksela nie zostanie
znaleziony żaden dopasowany komponent, tworzony jest nowy rozkład o~parametrach: $\hat{\pi}_{M+1} =
  \alpha$, $\hat{\vec{\mu}}_{M+1} = \vec{x}(t)$ oraz wariancji początkowej $\hat{\sigma}_{M+1}^2 =
  \sigma_0^2$~\cite{stauffer1999}. Jeśli liczba komponentów przekroczy dopuszczalne maksimum, rozkład
o~najmniejszej wadze jest usuwany~\cite{stauffer1999}.

\subsubsection{Ujemny rozkład a priori Dirichleta i eliminacja komponentów}

W~celu automatycznego określenia optymalnej liczby komponentów $M$ dla każdego piksela, Zivkovic
wprowadził do estymacji parametrów metodę maksymalizacji prawdopodobieństwa a posteriori (MAP)
z~użyciem sprzężonego rozkładu a priori Dirichleta~\cite{zivkovic2004, zivkovic2006}. Zdefiniowano
rozkład a priori w~postaci~\cite{zivkovic2004}:

\begin{equation}
  P = \prod_{m=1}^M \pi_m^{c_m}
  \label{eq:dirichlet_prior}
\end{equation}

Kluczowym elementem tej metody jest nadanie współczynnikom rozkładu a priori wartości ujemnych
($c_m = -c$)~\cite{zivkovic2004}. Z~matematycznego punktu widzenia ujemne współczynniki wprowadzają
karę za nadmierną złożoność modelu, co jest ściśle powiązane z~kryterium minimalnej długości opisu
(\textit{ang. Minimum Message Length} -- MML)~\cite{zivkovic2004}. Rozwiązanie równania MAP przy
uwzględnieniu ujemnego rozkładu a priori prowadzi do wyznaczenia wag
skorygowanych~\cite{zivkovic2004}:

\begin{equation}
  \hat{\pi}_m^{(t)} = \frac{1}{K} \left( \sum_{i=1}^t o_m^{(i)} - c \right)
  \label{eq:map_weights}
\end{equation}

gdzie parametr normalizacyjny wynosi $K = t - Mc$~\cite{zivkovic2004}. Aby zapobiec zanikaniu
wpływu kary wraz z~upływem czasu (gdy $t \to \infty$), ułamek $c/t$ zastępuje się stałą wartością
$c_T = c/T$~\cite{zivkovic2004}. Przy założeniu, że rzeczywista liczba komponentów opisujących
piksel jest niewielka, ostateczne równanie aktualizacji wag przyjmuje postać~\cite{zivkovic2004}:

\begin{equation}
  \hat{\pi}_m \leftarrow \hat{\pi}_m + \alpha (o_m^{(t)} - \hat{\pi}_m) - \alpha c_T
  \label{eq:mog2_final_weight_update}
\end{equation}

Zastosowanie stałej poprawki ujemnej $-\alpha c_T$ sprawia, że wagi komponentów, które nie są
regularnie wspierane przez nowe obserwacje, sukcesywnie maleją~\cite{zivkovic2004}. Gdy waga danego
komponentu osiągnie wartość ujemną, zostaje on trwale usunięty z~mieszaniny~\cite{zivkovic2004}.
W~ten sposób piksele reprezentujące stabilne elementy tła są opisywane zaledwie jednym rozkładem
Gaussa, natomiast piksele w~obszarach o~złożonej dynamice (np. migoczące światła, falująca
roślinność) utrzymują większą liczbę komponentów~\cite{zivkovic2004, zivkovic2006}.

\subsection{KNN \textit{(ang. K-nearest Neighbors)}}

Algorytm KNN w kontekście odejmowania tła stanowi alternatywne, nieparametryczne podejście do
modelowania rozkładu prawdopodobieństwa pikseli. W podejściu nieparametrycznym funkcja gęstości
jest szacowana bezpośrednio na podstawie danych historycznych, bez zakładania konkretnego kształtu
rozkładu~\cite{elgammal2000}. Podstawową techniką wykorzystywaną w tym podejściu jest estymacja
gęstości jądrowej (KDE, \textit{ang. Kernel Density Estimation}), która wymaga przechowywania
zbioru $N$ próbek historycznych $\vec{x}_n$ reprezentujących wcześniej obserwowane wartości
pikseli~\cite{elgammal2000}.

\subsubsection{Estymacja gęstości jądrowej z użyciem estymatora balonowego}

Podstawowy estymator Parzena–Rosenblatta szacuje gęstość prawdopodobieństwa punktu $\vec{x}$ w
$d$-wymiarowej przestrzeni cech na podstawie $N$ próbek historycznych $\vec{x}_n$ jako:

\begin{equation}
  \hat{p}(\vec{x}) = \frac{1}{N h^d} \sum_{n=1}^{N} K\!\left(\frac{\vec{x} - \vec{x}_n}{h}\right)
\end{equation}

gdzie $K(\cdot)$ jest funkcją jądrową, a $h$ stałą szerokością pasma. W kontekście odejmowania tła
szerokość pasma decyduje o wrażliwości klasyfikacji: jeżeli piksel znajduje się poza promieniem
wyznaczonym przez jądro, zostaje zaklasyfikowany jako pierwszy plan~\cite{zivkovic2006}. Użycie
stałej szerokości pasma może generować błędy w obszarach o zróżnicowanej gęstości danych —
nadmierne wygładzenie w regionach gęstych i niedoszacowanie w rzadkich~\cite{terrell1992}.

Zastosowanie estymatora balonowego (\textit{ang. balloon estimator})~\cite{terrell1992} rozwiązuje
ten problem poprzez uzależnienie szerokości jądra $h$ od lokalnej gęstości w punkcie
estymacji~$\vec{x}$. Szerokość okna dobiera się odwrotnie proporcjonalnie do lokalnej gęstości:

\begin{equation}
  h(\vec{x}) = \frac{k}{\left[N\,\hat{p}(\vec{x})\right]^{1/d}}
\end{equation}

Wówczas estymowana gęstość prawdopodobieństwa wyraża się wzorem:

\begin{equation}
  \hat{p}(\vec{x}) \approx \frac{1}{N \cdot h(\vec{x})^d}
  \sum_{n=1}^{N} K\!\left(\frac{\vec{x} - \vec{x}_n}{h(\vec{x})}\right)
\end{equation}

Dla optymalizacji wydajności obliczeniowej w systemach czasu rzeczywistego jako funkcję jądrową
$K(u)$ stosuje się jądro jednostajne (prostokątne)~\cite{terrell1992}. Wówczas estymacja gęstości
sprowadza się do zliczenia próbek $\vec{x}_n$ leżących wewnątrz sfery o promieniu $h(\vec{x})$
wyznaczonym przez odległość do $k$-tego sąsiada~\cite{zivkovic2006}:

\begin{equation}
  \hat{p}(\vec{x}) \approx \frac{k}{N \cdot V_d \cdot h(\vec{x})^d}
\end{equation}

gdzie $V_d$ oznacza objętość jednostkowej kuli w przestrzeni $d$-wymiarowej. Takie podejście
drastycznie skraca czas obliczeń w porównaniu z jądrem gaussowskim, przy zachowaniu adaptacyjności
szerokości pasma~\cite{zivkovic2006, mittal2004}.

\subsubsection{Metoda k-najbliższych sąsiadów w przestrzeni barw RGB}

W praktycznej implementacji algorytmu KNN klasyfikacja nowego piksela $\vec{x}(t)$ opiera się na
następującym schemacie decyzyjnym~\cite{zivkovic2006}:

\begin{enumerate}
  \item Zakłada się, że piksel jest reprezentowany jako wektor w przestrzeni barw RGB: $\vec{x}(t) =
          [R(t),\, G(t),\, B(t)]^T$.

  \item Oblicza się odległości euklidesowe $d_n$ między $\vec{x}(t)$ a każdą z $N$ próbek zgromadzonych w
        buforze historycznym tego piksela $\mathcal{X} = \{\vec{x}_1, \dots,
          \vec{x}_N\}$~\cite{zivkovic2006, cucchiara2003}.

  \item Zlicza się liczbę próbek $k^*$, dla których $d_n$ jest mniejsza od progu $D_{thr}$ (w OpenCV
        parametr \texttt{dist2Threshold} przechowuje kwadrat tej odległości):
        \begin{equation}
          k^* = \sum_{n=1}^{N} \mathbf{1}\!\left[d\!\left(\vec{x}(t),\,\vec{x}_n\right) < D_{thr}\right]
        \end{equation}
        gdzie $\mathbf{1}[\cdot]$ oznacza funkcję wskaźnikową. Próbki spełniające ten warunek
        mieszczą się wewnątrz sfery o promieniu $D_{thr}$ wokół $\vec{x}(t)$~\cite{prati2003,
          zivkovic2006, opencv2025}.

  \item Jeśli $k^* \ge k$, piksel zostaje zaklasyfikowany jako tło~\cite{zivkovic2006, opencv2025}. W
        przeciwnym razie jest oznaczany jako pierwszy plan~\cite{zivkovic2006, opencv2025}.
\end{enumerate}

Aktualizacja modelu w algorytmie KNN zachodzi w sposób ciągły — nowa próbka piksela zastępuje
najstarszą próbkę w buforze (schemat FIFO). Dzięki temu model stopniowo dostosowuje się do wolnych
zmian warunków oświetleniowych, takich jak ruch chmur czy dynamika cieni~\cite{zivkovic2006} .

\subsection{Ocena efektywności algorytmów}

Wyniki segmentacji pierwszego planu weryfikuje się przez porównanie wygenerowanej maski binarnej
z~maską referencyjną (\textit{ang. ground truth}) na poziomie pikseli~\cite{cdnet2014}. Każdy
piksel klasyfikowany jest jako jeden z~czterech przypadków: prawdziwy pozytyw (TP -- piksel
pierwszego planu wykryty poprawnie), fałszywy pozytyw (FP -- piksel tła błędnie oznaczony jako
pierwszy plan), prawdziwy negatyw (TN -- piksel tła wykryty poprawnie) oraz fałszywy negatyw (FN --
piksel pierwszego planu pominięty przez algorytm). Na tej podstawie wyznacza się następujące
metryki~\cite{cdnet2014}:

\begin{itemize}
  \item \textbf{Czułość} \textit{(ang. Recall / True Positive Rate)} — stosunek
        poprawnie wykrytych pikseli pierwszego planu do wszystkich pikseli pierwszego
        planu w masce referencyjnej:
        \begin{equation}
          \mathrm{Recall} = \frac{TP}{TP + FN}
        \end{equation}

  \item \textbf{Swoistość} \textit{(ang. Specificity / True Negative Rate)} — stosunek
        poprawnie sklasyfikowanych pikseli tła do wszystkich pikseli tła w~masce
        referencyjnej:
        \begin{equation}
          \mathrm{Specificity} = \frac{TN}{TN + FP}
        \end{equation}

  \item \textbf{Precyzja} \textit{(ang. Precision)} — stosunek poprawnie wykrytych
        pikseli pierwszego planu do wszystkich pikseli oznaczonych przez algorytm jako
        pierwszy plan:
        \begin{equation}
          \mathrm{Precision} = \frac{TP}{TP + FP}
        \end{equation}

  \item \textbf{Miara F1} \textit{(ang. F1-score)} — średnia harmoniczna precyzji
        i~czułości, zapewniająca zrównoważoną ocenę przy nierównomiernym rozkładzie
        klas:
        \begin{equation}
          F_1 = 2 \cdot \frac{\mathrm{Precision} \cdot \mathrm{Recall}}
          {\mathrm{Precision} + \mathrm{Recall}}
        \end{equation}

  \item \textbf{IoU} \textit{(ang. Intersection over Union)} — stosunek pola
        przecięcia maski wykrytej i~referencyjnej do pola ich sumy. W~kontekście
        ewaluacji BGS operuje się na maskach binarnych, nie na ramkach
        otaczających~\cite{cdnet2014}:
        \begin{equation}
          \mathrm{IoU} = \frac{TP}{TP + FP + FN}
        \end{equation}

  \item \textbf{PWC} \textit{(ang. Percentage of Wrong Classifications)} — odsetek
        błędnie sklasyfikowanych pikseli względem całkowitej liczby pikseli w~sekwencji;
        stosowany jako metryka syntetyczna w~benchmarku CDnet~\cite{cdnet2014}:
        \begin{equation}
          \mathrm{PWC} = 100 \cdot \frac{FP + FN}{TP + FP + TN + FN}
        \end{equation}
\end{itemize}

Wymienione metryki są standardowo stosowane w~benchmarku CDnet 2014~\cite{cdnet2014}, który stanowi
punkt odniesienia dla porównania algorytmów BGS w~niniejszej pracy. Metryki oceny śledzenia wielu
obiektów (MOTA, HOTA, IDF1) oraz reidentyfikacji (Rank-1/CMC) zostają omówione w~kolejnych
podrozdziałach poświęconych odpowiednio modułom detekcji i~śledzenia.

\subsection{Analiza porównawcza i ewaluacja efektywności}

\subsubsection{Porównanie wydajności algorytmów BGS w trudnych warunkach atmosferycznych}

Efektywność algorytmów odejmowania tła zależy od charakterystyki badanej sceny oraz rodzaju
zakłóceń środowiskowych~\cite{bouwmans2014}. Tabela~\ref{tab:iou_bgs_comparison} przedstawia
średnie wartości współczynnika IoU uzyskane przez wybrane algorytmy BGS w~badaniach porównawczych
dla różnego rodzaju zakłóceń atmosferycznych wideo~\cite{karbowiak2022}.

\begin{table}[htbp]
  \centering
  \caption{Średnie wartości współczynnika IoU wybranych algorytmów BGS. Żródło: opracowanie własne na podstawie danych z~\cite{karbowiak2022}.}
  \label{tab:iou_bgs_comparison}
  \begin{tabular}{|l|c|c|c|c|c|c|}
    \hline
    \textbf{Środowisko / Warunki}
            & \textbf{CNT} & \textbf{KNN} & \textbf{MOG} & \textbf{MOG2} & \textbf{GMG} & \textbf{SOM} \\ \hline
    Śnieg   & 0,41         & 0,38         & 0,44         & 0,25          & 0,42         & 0,52         \\ \hline
    Wiatr   & 0,13         & 0,13         & 0,33         & 0,11          & 0,12         & 0,29         \\ \hline
    Deszcz  & 0,17         & 0,14         & 0,37         & 0,10          & 0,31         & 0,25         \\ \hline
    Średnia & 0,28         & 0,32         & 0,39         & 0,25          & 0,42         & 0,47         \\ \hline
  \end{tabular}
\end{table}

Analiza wyników przedstawionych w~Tabeli~\ref{tab:iou_bgs_comparison} pozwala na sformułowanie
kilku obserwacji dotyczących zachowania algorytmów BGS w~trudnych warunkach atmosferycznych.

Najsłabszą i~najbardziej równomierną odporność na zakłócenia środowiskowe wykazują algorytmy KNN
oraz MOG2 -- oba uzyskują najniższe wartości IoU zarówno przy wietrze (0,13 i~0,11), jak i~przy
deszczu (0,14 i~0,10). Algorytm MOG z~kolei wykazuje zaskakująco wysoką odporność w~warunkach
wietrznych (IoU~=~0,33) oraz deszczowych (IoU~=~0,37), co wynika z~jego większej bezwładności
modelu -- stała liczba komponentów gaussowskich wolniej absorbuje szybkie
zakłócenia~\cite{stauffer1999}.

Spośród wszystkich badanych metod najwyższą średnią skuteczność osiąga SOM (IoU~=~0,47),
a~najniższą -- MOG2 (IoU~=~0,25). Algorytm SOM (\textit{ang. Self-Organizing Map}) modeluje tło za
pomocą samoorganizacji poprzez szczutczną sieć neuronową\cite{karbowiak2022}. Dzięki temu
skuteczniej radzi sobie ze złożonymi problemami, takimi jak ruszające się tło czy zmiany w
oświetleniu~\cite{karbowiak2022}.

Algorytm GMG osiąga wysoką średnią (IoU~=~0,42) dzięki bayesowskiemu modelowi tła budowanemu na
podstawie pierwszych klatek sekwencji (domyślnie 120 klatek w~OpenCV). Wysoka skuteczność przy
śniegu (IoU~=~0,42) jest efektem dobrego odwzorowania statystycznej zmienności tła w~fazie
inicjalizacji. Metoda ta ma jednak istotne ograniczenie -- wymaga stabilnej fazy uczenia i~słabo
radzi sobie z~dynamicznymi scenami od pierwszej klatki~\cite{karbowiak2022}.

Algorytm CNT (\textit{ang. Counter}, \textit{ang. Count of Non-stationary Targets}) opiera się na
zliczaniu, ile razy wartość piksela zmieniła się w danym czasie. Jest to metoda prosta
obliczeniowo, osiągająca zadowalające wyniki dla śniegu (IoU~=~0,41), jest zawodna przy wietrze
i~deszczu (IoU~=~0,13 i~0,17).

Pomimo że algorytmy SOM, GMG i~CNT uzyskują wyższe wartości IoU od KNN i~MOG2, w~niniejszej pracy
zdecydowano się na implementację tych ostatnich. Wynika to z~ich natywnej dostępności i~pełnego
wsparcia w~bibliotece OpenCV, niskiego zapotrzebowania na zasoby obliczeniowe niezbędnego
w~przetwarzaniu strumienia wideo w~czasie rzeczywistym, a~także szczegółowej dokumentacji
i~ugruntowanej pozycji w~literaturze~\cite{zivkovic2006}. Algorytmy SOM i~GMG wymagają odpowiednio
dłuższej fazy inicjalizacji lub wyższych zasobów, co utrudnia ich zastosowanie w~scenariuszach
z~dynamicznie zmieniającym się tłem bez etapu wstępnego uczenia.

\subsubsection{Wydajność obliczeniowa, ograniczenia sprzętowe i zjawisko dryfu tła}

W systemach czasu rzeczywistego wydajność obliczeniowa jest istotnym czynnikiem przy wdrażaniu
oprogramowania\cite{bouwmans2014}. MOG2, dzięki dynamicznemu doborowi liczby komponentów
gaussowskich opartemu na kryterium MML \textit{(ang. Minimum Message Length)}~\cite{zivkovic2004},
charakteryzuje się niskim zużyciem pamięci RAM. Pozwala to na jego stosowanie w~architekturach
wielokamerowych na platformach wbudowanych~\cite{waghmare2026}.

Najpoważniejszą wadą klasycznych algorytmów BGS, w~tym MOG2 i~KNN, jest podatność na zjawisko dryfu
modelu tła (\textit{ang. background drift})~\cite{waghmare2026}. Jeżeli obiekt pierwszego plany
pozostanie nieruchomy, w ciągu 100 klatek może zostać zaklasyfikowany jako część
tła~\cite{waghmare2026}. Skutkuje to fragmentacją ścieżki śledzenia i~utratą ciągłości detekcji.
Dlatego na trudnych zbiorach testowych samodzielne algorytmy BGS osiągają niską czułość
detekcji~\cite{waghmare2026}.

%%%%%%%%%%%%%%%%%%%%%%%%%%%%%%%%%%%%%%%%%%%%%%%%%%%%%%%%%%%
%%%%%%%%%%%%%%%%%%%%%%%%%%%%%%%%%%%%%%%%%%%%%%%%%%%%%%%%%%%

\section{Detekcja obiektów oparta na głębokim uczeniu (Deep Learning)}

Przed upowszechnieniem wielowarstwowych sieci neuronowych, systemy detekcji obiektów opierały się
na ręcznie projektowanych cechach wizualnych \textit{(ang. hand-crafted features)}, które były
następnie klasyfikowane za pomocą tradycyjnych algorytmów uczenia maszynowego. Jednym
z~najważniejszych rozwiązań w~obszarze detekcji obiektów w~czasie rzeczywistym stał się algorytm
Viola-Jones, pierwotnie opracowany do detekcji twarzy. Metoda ta opiera się na trzech kluczowych
koncepcjach: reprezentacji obrazu całkowego \textit{(ang. Integral Image)}, wyznaczaniu prostych
cech \textit{(ang. Haar-like features)} oraz kaskadowej strukturze klasyfikatorów opartej na
algorytmie AdaBoost \cite{viola2001}. Zastosowanie obrazu całkowego umożliwiło obliczanie sum
jasności pikseli wewnątrz prostokątnych obszarów w czasie stałym $O(1)$, niezależnie od skali.
Pomimo rewolucyjnej szybkości działania, detektor Viola-Jones wykazywał niską odporność na rotacje
obiektów w przestrzeni trójwymiarowej, okluzje oraz gwałtowne zmiany kontrastu i oświetlenia sceny
\cite{viola2001}. Kolejny przełomowy algorytm, służący do detekcji sylwetek ludzkich, to metoda
HOG~+~SVM. Algorytm opierał się na histogramach zorientowanych gradientów liczonych w komórkach
obrazu i normalizowanych lokalnie w nakładających się blokach, co poprawiało odporność na zmiany
kontrastu \cite{dalal2005}. Autorzy wykazali, że taka reprezentacja znacznie przewyższa
wcześniejsze zestawy cech w zadaniu detekcji ludzi. Ograniczeniem podejścia pozostawała m.in.
podatność na deformacje postawy oraz koszt obliczeniowy wynikający ze skanowania obrazu metodą
przesuwanego okna \cite{dalal2005}.

\subsection{Splotowe sieci neuronowe (CNN) i detektory dwuetapowe}
Rok 2012 przyniósł zasadnicze zmiany w wizji komputerowej wraz z sukcesem głębokiej sieci splotowej
AlexNet w konkursie ILSVRC \cite{krizhevsky2012imagenet}. Model ten, zbudowany z pięciu warstw
splotowych i trzech warstw w pełni połączonych, pokazał, że sieci neuronowe mogą samodzielnie uczyć
się użytecznych cech obrazu lepiej niż metody ręcznie projektowane \cite{krizhevsky2012imagenet}.
Zastosowanie funkcji ReLU oraz technik regularizacji, takich jak dropout, umożliwiło skuteczne
trenowanie dużych modeli na procesorach graficznych (GPU).

Zastosowanie sieci splotowych do lokalizacji obiektów doprowadziło do powstania detektorów
dwuetapowych. W pierwszym kroku system wskazywał możliwe miejsca występowania obiektu, a w drugim
klasyfikował je i doprecyzowywał położenie. Pierwszym istotnym modelem był R-CNN, który łączył
tradycyjne generowanie propozycji regionów, na przykład metodą Selective Search, z siecią CNN
analizującą każdy region osobno \cite{girshick2014rich}. Było to jednak wolne, ponieważ ten sam
fragment obrazu był przetwarzany wiele razy. Fast R-CNN przyspieszył ten proces, obliczając cechy
dla całego obrazu tylko raz i wykorzystując warstwę ROI Pooling do pobierania informacji dla
proponowanych regionów \cite{girshick2015fast}. Dalszym ulepszeniem był Faster R-CNN, który
zastąpił Selective Search siecią RPN, czyli modułem proponującym regiony bezpośrednio na podstawie
map cech \cite{ren2015faster}. Dzięki temu model stał się znacznie szybszy i bardzo dokładny, choć
nie zawsze osiągał pełny czas rzeczywisty.

\subsection{Detektory jednoetapowe}

Modele dwuetapowe, mimo wysokiej dokładności, często były zbyt wolne do zastosowań wymagających
szybkiej reakcji, takich jak monitoring czasu rzeczywistego czy systemy wbudowane. Odpowiadając na
tę potrzebę,powstały detektory jednoetapowe (ang. \textit{one-stage detectors}), które bezpośrednio
przekształcają obraz wejściowy w predykcje klas oraz współrzędnych ramek otaczających, bez osobnego
etapu generowania propozycji regionów.

Jednym z najważniejszych algorytmów tej grupy stał się YOLO \textit{(ang. You Only Look Once)}, w
którym potraktowano detekcję obiektów jako pojedynczy problem regresji realizowany przez jedną sieć
neuronową w jednym przebiegu obrazu (\textit{single forward pass}) \cite{redmon2016yolo}. Dzięki
temu możliwe stało się wykrywanie obiektów w czasie zbliżonym do rzeczywistego, przy zachowaniu
dobrej skuteczności. Kolejne wersje tej rodziny rozwijały to podejście, poprawiając kompromis
między szybkością a dokładnością. YOLOv7 wprowadził architekturę E-ELAN \textit{(ang. Extended
  Efficient Layer Aggregation Network)}, która usprawniła przepływ informacji w sieci, oraz koncepcję
\textit{compound scaling}, pozwalającą lepiej dobierać rozmiar modelu do możliwości obliczeniowych.
Autorzy zastosowali także tzw. \textit{bag-of-freebies}, czyli zestaw ulepszeń poprawiających
jakość uczenia bez zwiększania kosztu inferencji \cite{wang2023yolov7}. YOLOv8 przeszedł na
podejście bezkotwicowe (\textit{anchor-free}), rezygnując z wcześniej definiowanych ramek bazowych,
oraz zastosował rozdzieloną głowę detekcyjną \textit{(ang. split head)}, w której klasyfikacja i
regresja są wykonywane osobno \cite{ultralytics2023yolov8}. Z kolei YOLOv9 zaproponował rozwiązania
mające ograniczać utratę informacji w głębokich sieciach oraz poprawiać przepływ gradientu, co ma
szczególne znaczenie w przypadku małych i trudnych do wykrycia obiektów. \cite{wang2024yolov9}.

Alternatywnym podejściem do detekcji jednoetapowej jest SSD \textit{(ang. Single Shot MultiBox
  Detector)}. Metoda ta wykorzystuje kilka map cech o różnych rozdzielczościach, dzięki czemu lepiej
radzi sobie z obiektami o zróżnicowanych rozmiarach \cite{liu2016ssd}. W porównaniu z modelami
dwuetapowymi SSD oferował prostszą architekturę i większą szybkość działania, choć początkowo
często ustępował im dokładnością.

Głównym problemem detektorów jednoetapowych była jednak silna nierównowaga klasowa: w obrazie
dominuje tło, a obiekty zajmują tylko niewielką część danych. Prowadziło to do sytuacji, w której
model był zbyt mocno uczony rozpoznawania tła, a zbyt słabo skupiał się na rzeczywistych obiektach
\cite{lin2017focal}. Problem ten skutecznie złagodził RetinaNet, który wprowadził funkcję straty
Focal Loss:
\[
  FL(p_t)=-\alpha_t(1-p_t)^\gamma\log(p_t)
\]
Czynnik modulujący $(1 - p_t)^\gamma$ (gdzie $\gamma$ to parametr skupienia) sprawia, że w tym
ujęciu łatwe przykłady, czyli takie, które model klasyfikuje z dużą pewnością, otrzymują mniejszą
wagę w procesie uczenia. Dzięki temu gradient jest bardziej skoncentrowany na trudnych przykładach
i na obiektach pierwszoplanowych, co poprawia skuteczność detekcji bez rezygnacji z szybkości
charakterystycznej dla podejścia jednoetapowego \cite{lin2017focal}.

\subsection{Transformery w detekcji obiektów}

Wprowadzenie mechanizmów uwagi, znanych wcześniej głównie z przetwarzania języka naturalnego,
doprowadziło do powstania architektury DETR (\textit{Detection Transformer}). Model ten traktuje
detekcję obiektów jako problem bezpośredniego przewidywania zbioru wyników (\textit{direct set
  prediction}). W praktyce oznacza to, że zamiast stosować kotwice (\textit{anchors}) i dodatkowy
etap usuwania nadmiarowych ramek NMS \textit{(ang. Non-Maximum Suppression)}, DETR łączy ekstraktor
cech oparty na sieci splotowej z enkoderem i dekoderem Transformer oraz mechanizmem dopasowania
dwudzielnego opartym na algorytmie węgierskim \cite{carion2020end}.

Mimo eleganckiej konstrukcji oryginalny DETR miał istotne ograniczenia. Model uczył się bardzo
wolno, wymagał dużych zasobów obliczeniowych i osiągał słabsze wyniki w detekcji małych obiektów.
Odpowiedzią na te problemy był Deformable DETR. Zastosowano w nim odkształcalną uwagę
(\textit{deformable attention}), która zamiast analizować cały obraz, skupia się na niewielkiej
liczbie punktów wybranych wokół punktów referencyjnych. Dzięki temu model lepiej wykorzystuje cechy
wieloskalowe, szybciej się uczy i radzi sobie skuteczniej z obiektami o różnych rozmiarach
\cite{zhu2021deformable}.

Nowszym rozwinięciem tej linii jest RT-DETR (\textit{Real-Time Detection Transformer}),
zaprojektowany z myślą o zastosowaniach czasu rzeczywistego. Model ten łączy wydajny enkoder
hybrydowy z mechanizmem wyboru najbardziej obiecujących zapytań dekodera, co pozwala ograniczyć
koszty obliczeniowe i zachować wysoką szybkość działania. W odróżnieniu od klasycznych detektorów
opartych na CNN RT-DETR dąży do osiągnięcia bardzo dobrego kompromisu między dokładnością a czasem
inferencji, a przy tym eliminuje potrzebę stosowania NMS jako osobnego etapu przetwarzania
\cite{zhao2024detrs}.

\subsection{Dane treningowe a problem luki domenowej}

Skuteczność głębokich modeli detekcji zależy od wielkości i jakości danych uczących. Modele są
zazwyczaj trenowane na dużych, ogólnych zbiorach badawczych, takich jak MS COCO lub ImageNet.
Jednak bezpośrednie przeniesienie modelu wytrenowanego na danych syntetycznych do rzeczywistego
środowiska przemysłowego lub systemu CCTV zawodzi, ujawniając problem luki domenowej \textit{(ang.
  domain shift)}\cite{he2024}.

Luka domenowa wynika z różnic w charakterystyce obrazów między zbiorem treningowym a rzeczywistym
środowiskiem.
\begin{itemize}
  \item Zdjęcia w zbiorach akademickich najczęściej wykonywane są z perspektywy na wysokości oczu, kamery
        monitoringu są zwykle umieszczone znacznie wyżej. Powoduje to zniekształcenia perspektywy, które
        wpływają na kształt i rozmiar obiektów, a także na ich wzajemne relacje przestrzenne.
  \item Kamery CCTV charakteryzują się specyficznymi zniekształceniami optycznymi (np. obiektywy
        szerokokątne typu rybie oko), silną kompresją stratną strumienia wideo oraz wysokim poziomem szumów
        matrycy w warunkach nocnych.
  \item Zbiory testowe rzadko reprezentują skrajne zjawiska pogodowe (silny deszcz, śnieżyca, mgła,
        oślepiające słońce generujące głębokie cienie), które w rzeczywistym monitoringu zewnętrznym
        występują powszechnie \cite{he2024}.
\end{itemize}

Najczęściej stosowane techniki niwelowania luki domenowej w detekcji obiektów można podzielić na
trzy główne grupy: translację domen, samouczenie z pseudoetykietami oraz dopasowanie cech między
domenami. Pierwsza grupa polega na przekształcaniu danych źródłowych tak, aby bardziej przypominały
dane docelowe, zwykle przy użyciu modeli generatywnych. Wiąże się to jednak z utratą informacji i
wysoką złożonością obliczeniową. \cite{he2024}.

Metody samouczenia \textit{(ang. self-training)}, w których model uczony jest na danych źródłowych
z etykietami oraz na danych docelowych opisanych pseudoetykietami generowanymi przez model
nauczyciela. Jego skuteczność zależy jednak od jakości pseudoetykiet, dlatego wiele nowszych metod
wprowadza mechanizmy ich selekcji, ważenia lub filtrowania \cite{he2024}.

Szeroko stosowane metody dopasowania cech (\textit{domain alignment}), dążą do zmniejszenia różnic
pomiędzy rozkładami cech w domenie źródłowej i docelowej. W klasycznych rozwiązaniach wykorzystuje
się do tego sieci przeciwstawne, w których dyskryminator próbuje odróżnić cechy pochodzące z obu
domen, a ekstraktor cech uczy się tworzyć reprezentacje nieodróżnialne z punktu widzenia domeny. W
detekcji obiektów takie dopasowanie może zachodzić na poziomie obrazu, regionu lub pojedynczej
instancji \cite{chen2018, he2024}.

Nowsze badania wskazują jednak, że globalne wyrównywanie cech na poziomie całego obrazu szkodzi
precyzji lokalizacji detektorów, gdyż prowadzi do transferu informacji z nieistotnych obszarów tła
\cite{he2024}. W związku z tym wprowadzono zaawansowane strategie adaptacji różnicowej (np. UFOA –
Uncertainty-based Foreground-Oriented Alignment) \cite{he2024}, które selektywnie nadają wyższy
priorytet obszarom reprezentującym obiekty pierwszoplanowe.

\subsection{Detekcja w trudnych warunkach}

Wdrożenie skutecznych algorytmów detekcji obiektów w systemach monitoringu wymaga wysokiej
odporności modeli na zakłócenia. Do najczęstszych problemów należą okluzje, gwałtowne zmiany
oświetlenia oraz niska rozdzielczość małych lub odległych obiektów \cite{nikouei2025, cheng2026}.

Techniki ograniczające wpływ okluzji obejmują między innymi:
\begin{itemize}
  \item \textbf{Generatywne warstwy kompozycyjne} – w tym podejściu klasyczne głowice klasyfikacyjne zastępowane są strukturą generatywną, która potrafi wnioskować o niewidocznych częściach obiektu na podstawie relacji przestrzennych między jego widocznymi fragmentami. Przykładem takiego rozwiązania jest model CompositionalNet \cite{cheng2026}.
  \item \textbf{Algorytmy głosowania częściowego \textit{(ang. part-based voting)}} – model uczy się niezależnych reprezentacji poszczególnych części obiektu, a następnie szacuje położenie całego obiektu na podstawie wykrytych elementów \cite{cheng2026}. W detekcji pieszych może to oznaczać osobne wykrywanie głowy, ramion i nóg, a następnie wyznaczenie położenia całej sylwetki.
  \item \textbf{Łączenie wielu widoków i wspólne wnioskowanie} – w systemach wielokamerowych obiekt zasłonięty z jednej perspektywy może być dobrze widoczny z innej. Połączenie predykcji z kilku sąsiadujących kamer pozwala trafniej odtworzyć położenie obiektu, nawet jeśli w pojedynczym obrazie jest on silnie zasłonięty \cite{cheng2026}.
\end{itemize}

Systemy CCTV działają przez całą dobę, dlatego są narażone na zmienne warunki oświetleniowe
utrudniające detekcję. Problem stanowi zarówno intensywne światło słoneczne powodujące głębokie
cienie, jak i nagłe prześwietlenia, na przykład od reflektorów samochodowych w nocy, a także bardzo
słabe oświetlenie po zmroku \cite{nikouei2025}. Dodatkowo opady atmosferyczne, takie jak deszcz i
śnieg, wprowadzają zakłócenia do obrazu i osłabiają składowe wysokoczęstotliwościowe, co prowadzi
do rozmycia krawędzi obiektów \cite{wille2026}. Jednym ze sposobów ograniczania utraty detali są
moduły o zmiennym polu recepcyjnym, takie jak RepHELAN, które wykorzystują heterogeniczne
konwolucje do zwiększania efektywnego pola recepcyjnego bez pogorszenia reprezentacji małych
obiektów \cite{yang2024}. Z kolei w trudnych warunkach oświetleniowych stosuje się mechanizmy uwagi
wieloskalowej, na przykład LSKA (\textit{Large Separable Kernel Attention}), które mogą
adaptacyjnie ograniczać wpływ lokalnych prześwietleń i koncentrować się na bardziej stabilnych
zależnościach przestrzennych \cite{tang2026}.

W rozległych scenach monitoringu, takich jak place miejskie, skrzyżowania czy parkingi, obiekty
znajdujące się daleko od kamery klasyfikowane są jako małe jeżeli zajmują na obraziemniej niż $32
  \times 32$, a~często stanowią nawet mniej niż $20 \times 20$ pikseli \cite{nikouei2025}.
W~klasycznych sieciach konwolucyjnych kolejne operacje zmniejszania rozdzielczości, realizowane na
przykład przez konwolucje ze skokiem lub warstwy MaxPooling, mogą prowadzić do utraty informacji o
tak małych obiektach jeszcze przed przekazaniem cech do głowic predykcyjnych \cite{nikouei2025}.
Aby temu przeciwdziałać, nowoczesne architektury wykorzystują mechanizmy wzmacniające sygnał małych
obiektów. Jednym z nich są gęste połączenia wspomagające, takie jak moduł SAF (\textit{Superficial
  Assisted Fusion}) stosowany w części agregującej cechy (\textit{neck}), który przekazuje do
głębszych warstw informacje o wysokiej rozdzielczości przestrzennej z początkowych etapów
ekstrakcji cech \cite{yang2024}. Innym rozwiązaniem jest regresja dystrybucyjna, na przykład DFL
(\textit{Distribution Focal Loss}) wykorzystywana w modelach YOLOv8 i YOLOv9, w której współrzędne
ramki nie są modelowane jako pojedyncze wartości, lecz jako rozkłady prawdopodobieństwa wokół jej
krawędzi. Pozwala to precyzyjniej opisać niepewność lokalizacji małego obiektu, także na poziomie
subpikselowym \cite{yang2024}.
%%%%%%%%%%%%%%%%%%%%%%%%%%%%%%%%%%%%%%%%%%%%%%%%%%%%%%%%%%%
%%%%%%%%%%%%%%%%%%%%%%%%%%%%%%%%%%%%%%%%%%%%%%%%%%%%%%%%%%%

\section{Śledzenie wielu obiektów w pojedynczej kamerze}

Śledzenie wielu obiektów (Multiple Object Tracking – MOT) w strumieniu wideo z pojedynczej kamery polega na jednoczesnym szacowaniu pozycji wielu poruszających się obiektów w czasie oraz na zachowaniu ich spójnej tożsamości w kolejnych klatkach sekwencji obrazów. Zadanie to wymaga rozwiązania trzech podstawowych problemów:
\begin{itemize}
  \item \textbf{Asocjacji danych} \textit{(ang. data association)} - powiązania niezależnych detekcji tych samych obiektów na przestrzeni kolejnych klatek wideo i przypisania im stałego, unikalnego identyfikatora ang. \textit{(ang. track ID)} \cite{bewley2016}.
  \item \textbf{Inicjalizacji nowych ścieżek} — wykrywania momentu wejścia nowego, dotąd nieśledzonego obiektu w pole widzenia kamery i utworzenia dla niego nowej trajektorii \cite{bewley2016}.
  \item \textbf{Terminacji ścieżek} — identyfikacji momentu, w którym obiekt opuszcza monitorowaną scenę (lub zostaje długotrwale zasłonięty), w celu zakończenia aktualizacji jego trajektorii. Operacja ta zapobiega nieograniczonemu wzrostowi liczby śledzonych obiektów oraz redukuje błędy lokalizacji wynikające z długotrwałej predykcji ruchu \cite{bewley2016}.
\end{itemize}
\subsection{Paradygmat śledzenia ścieżek \textit{(ang. tracking-by-detection)}}
Od około 2015 roku w dziedzinie wieloobiektowego śledzenia dominuje paradygmat śledzenia przez detekcję \textit{(ang. tracking-by-detection lub detection-based tracking)} – DBT \cite{wojke2017}. Podejście to opiera się na rozdzieleniu modułu lokalizacji przestrzennej od modułu odpowiedzialnego za utrzymanie tożsamości. Przepływ informacji realizowany jest kaskadowo: najpierw algorytm detekcyjny (np. głęboka sieć neuronowa z rodziny YOLO) niezależnie lokalizuje obiekty w każdej klatce wideo, a następnie moduł śledzący łączy te dyskretne predykcje w ciągłe trajektorie \cite{bewley2016, wojke2017}. Taka separacja w architekturze systemu sprawia, że precyzja i czułość całego systemu MOT są bezpośrednio ograniczone przez jakość działania zastosowanego detektora bazowego \cite{bewley2016}.

\subsection{Metody predykcji stanu i algorytmy skojarzenia}
Aby efektywnie połączyć detekcje z klatki $t$ z obiektami w klatce $t+1$, nowoczesne systemy
śledzenia wykorzystują kombinację filtrów estymacji stanu oraz algorytmów optymalizacji dyskretnej.

\subsubsection{Filtr Kalmana}
Do modelowania dynamiki ruchu obiektów stosuje się liniowy filtr Kalmana \cite{kalman1960,
  bewley2016}. Zakładając model ruchu o stałej prędkości na płaszczyźnie obrazu, filtr ten przewiduje
przyszłą pozycję i wymiary ramki ograniczającej obiekt w kolejnej klatce wideo. Pozwala to zawęzić
obszar poszukiwań i stabilizuje proces śledzenia nawet w przypadku chwilowych zakłóceń
\cite{bewley2016}.

\subsubsection{Algorytm węgierski}
Powiązanie przewidzianych pozycji z nowymi detekcjami formułuje się jako problem skojarzenia w
grafie dwudzielnym (\textit{ang. bipartite matching}) o minimalnym koszcie \cite{bewley2016}.
Optymalne rozwiązanie tego zadania zapewnia algorytm węgierski \cite{kuhn1955, bewley2016}. Metoda
ta szuka optymalnych par w relacji jeden do jednego dla przewidzianych trajektorii i nowych
detekcji, tak aby zminimalizować łączny koszt – np. odległość opartą na stopniu nakładania się
ramek (IoU) \cite{kuhn1955, bewley2016}.

\subsection{Ewolucja kluczowych algorytmów śledzenia}
Rozwój algorytmów klasy MOT doprowadził do powstania kilku przełomowych architektur, stopniowo
rozwiązujących problemy okluzji i wydajności obliczeniowej.

\subsubsection{SORT \textit(Simple Online and Realtime Tracking), 2016:}
Minimalistyczny algorytm łączący predykcję filtrem Kalmana i dopasowanie oparte na
metryce IoU za pomocą algorytmu Węgierskiego. Dzięki prostej budowie cechuje się dużą szybkością, aktualizuje ścieżki z częstotliwościom 260 $Hz$, lecz był podatny na natychmiastowe gubienie tożsamości w przypadku nawet krótkotrwałego
zasłonięcia obiektu

\subsubsection{DeepSORT, 2017}
Rozwinięcie SORT o głęboki model asocjacji wyglądu. Poprzez zastosowanie dedykowanej sieci CNN
wytrenowanej na zadanie reidentyfikacji (Re-ID), DeepSORT wyznacza deskryptory wizualne obiektów,
co pozwala na poprawne powiązanie tożsamości po długotrwałych okluzjach, drastycznie redukując
liczbę przełączeń tożsamości.

\subsubsection{ByteTrack, 2022}
Nowatorski algorytm zwracający uwagę na błąd systematyczny polegający na odrzucaniu detekcji o
niskim poziomie ufności (poniżej progu detekcji, np. $0.5$), które często reprezentują obiekty
częściowo zasłonięte lub rozmyte. ByteTrack wprowadził dwustopniowy mechanizm asocjacji BYTE, który
najpierw dopasowuje detekcje o wysokiej ufności, a w drugim kroku kojarzy detekcje o niskiej
ufności z pozostałymi, nieprzypisanymi trajektoriami przy użyciu podobieństwa IoU, co znacząco
ogranicza fragmentację ścieżek.

\subsubsection{BoT-SORT, 2022}
Zestaw optymalizacji (ang. bag of tricks) dla ByteTrack, wprowadzający:

Kompensację ruchu kamery (CMC): Poprzez estymację globalnego przepływu tła, co zapobiega
rozbieżnościom filtru Kalmana przy drganiach kamery.

Zmodyfikowany wektor stanu: Bezpośrednie modelowanie szerokości i wysokości ramki zamiast jej
proporcji, co poprawiło precyzję dopasowania predykcji.

Zintegrowaną macierz kosztów: Łączącą odległość IoU i podobieństwo cech wyglądu (Re-ID) w sposób
komplementarny.

\subsection{Ocena efektywności algorytmów MOT}
skopiować z wczesnijeszego + benchmarki

\subsection{Wyzwania i ograniczenia w praktycznych zastosowaniach}

Aplikacja algorytmów śledzenia w rzeczywistych sieciach kamer monitoringu wizyjnego napotyka na
bariery fizyczne i środowiskowe:
\begin{itemize}
  \item \textbf{Okluzje} — Chwilowe lub długotrwałe zasłonięcie obiektu przez elementy infrastruktury miejskiej (latarnie, drzewa) powoduje dryf predykcji filtru Kalmana i w konsekwencji fragmentację trajektorii lub zmianę identyfikatora.
  \item \textbf{Wysoka gęstość obiektów} — Powoduje nakładanie się sylwetek, co degraduje jakość ekstrakcji cech wizualnych i uniemożliwia poprawną pracę klasycznych modeli Re-ID z powodu podobieństwa cech otoczenia.
  \item \textbf{Szybki i nieliniowy ruch} — Narusza założenia o jednostajności ruchu stosowane w filtrach Kalmana, prowadząc do niepoprawnego pozycjonowania ramek predykcyjnych w klatce $t+1$ i błędów asocjacji.
\end{itemize}

%%%%%%%%%%%%%%%%%%%%%%%%%%%%%%%%%%%%%%%%%%%%%%%%%%%%%%%%%%%
%%%%%%%%%%%%%%%%%%%%%%%%%%%%%%%%%%%%%%%%%%%%%%%%%%%%%%%%%%%

\section{Reidentyfikacja obiektów (Re-ID)}
\label{sec:reidentyfikacja_reid}

\subsection{Definicja i sformułowanie problemu reidentyfikacji (Re-ID)}
\label{subsec:definicja_reid}

Reidentyfikacja obiektów (\textit{ang. person/vehicle re-identification} -- Re-ID) stanowi kluczowy
komponent nowoczesnych, rozproszonych systemów nadzoru wizyjnego, umożliwiający kojarzenie
tożsamości celów ruchomych przemieszczających się pomiędzy fizycznie rozłącznymi obszarami
obserwowanymi przez kamery o~niepokrywających się polach widzenia (\textit{ang. non-overlapping
  fields of view})~\cite{ristani2016, zheng2016person}. Z~formalnego punktu widzenia, proces ten
definiuje się jako zadanie wyszukiwania informacji wizualnej (\textit{ang. image
  retrieval})~\cite{zheng2016person}. Dane wejściowe dla algorytmu Re-ID składają się z~dwóch
podstawowych zbiorów:
\begin{itemize}
  \item \textbf{Zapytania} (\textit{ang. query} lub \textit{probe}) -- reprezentowanego przez wycięty kadr obrazu zawierający sylwetkę poszukiwanego obiektu (człowieka lub pojazdu), najczęściej wygenerowany przez detektor lub moduł śledzenia jedno-kamerowego~\cite{ristani2016, zheng2016person},
  \item \textbf{Galerii} (\textit{ang. gallery}) -- bazy danych zawierającej kadry obiektów referencyjnych zarejestrowanych przez inne kamery w~sieci, w~której poszukuje się dopasowań~\cite{zheng2016person}.
\end{itemize}
Zadanie algorytmu polega na ekstrakcji reprezentacji cech (\textit{ang. feature representation}) z~zapytania oraz elementów galerii, a~następnie na obliczeniu miary podobieństwa semantycznego (odległości w~przestrzeni cech) i~uszeregowaniu (rangowaniu) bazy galerii od najbardziej do najmniej podobnego kandydata względem zapytania~\cite{zheng2016person}. Głównym wyzwaniem w~tak sformułowanym problemie jest konieczność opierania się wyłącznie na cechach wyglądu zewnętrznego (np. kolorze i~strukturze odzieży), ponieważ niska rozdzielczość sensorów CCTV oraz ukośne kąty zawieszenia kamer uniemożliwiają ekstrakcję jednoznacznych cech biometrycznych, takich jak geometria twarzy czy tekstura tęczówki oka~\cite{ristani2016, zheng2016person}.

\subsection{Klasyczne podejścia oparte na cechach projektowanych ręcznie}
\label{subsec:hand_crafted_reid}

Przed przełomem w~dziedzinie głębokiego uczenia, badania nad reidentyfikacją osób koncentrowały się
na ręcznym projektowaniu deskryptorów wizualnych (\textit{ang. hand-crafted features}) odpornych na
transformacje geometryczne i~fotometryczne sceny oraz na optymalizacji klasycznych miar odległości
(\textit{ang. metric learning})~\cite{zheng2016person}. Do pionierskich i~najczęściej cytowanych
rozwiązań z~tej ery należą:
\begin{itemize}
  \item \textbf{Histogramy koloru}: Reprezentujące globalny lub lokalny rozkład barw w~różnych przestrzeniach (RGB, HSV, YUV). Metody te charakteryzowały się prostotą obliczeniową, lecz wykazywały skrajną wrażliwość na zmiany oświetlenia i~efekty metameryzmu.
  \item \textbf{SDALF} (\textit{ang. Symmetry-Driven Accumulation of Local Features}): Algorytm zaproponowany przez Farenzenę i~in.~\cite{farenzena2010person}. SDALF opiera się na lokalizacji osi symetrii i~asymetrii ludzkiego ciała, co pozwala na podział sylwetki na stabilne regiony (głowa, tułów, nogi). Cechy lokalne (takie jak tekstury oparte na filtrach Gabora i~histogramy kolorów) są akumulowane wzdłuż tych osi, co pozwala na uzyskanie reprezentacji odpornej na rotacje obiektu i~zmiany pozy~\cite{farenzena2010person}.
  \item \textbf{LOMO} (\textit{ang. Local Maximal Occurrence}): Metoda opracowana przez Liao i~in.~\cite{liao2015person}. Deskryptor LOMO analizuje poziome pasy obrazu, wyznaczając lokalne histogramy kolorów (na bazie transformaty Retinexa neutralizującej zmiany oświetlenia) oraz tekstur (za pomocą skali invariantnej operatora SILTP)~\cite{liao2015person}. Maksymalizacja wartości wystąpień cech wzdłuż linii poziomych drastycznie zwiększa odporność modelu na zmiany punktu widzenia kamery (\textit{ang. viewpoint invariance})~\cite{liao2015person}.
\end{itemize}

\subsection{Uczenie reprezentacji metrycznych i sieci siamskie}
\label{subsec:metric_learning_reid}

Zastosowanie głębokich sieci neuronowych zrewolucjonizowało Re-ID poprzez zjednoczenie etapu
ekstrakcji cech z~uczeniem optymalnej przestrzeni embeddingów, gdzie odległość euklidesowa
bezpośrednio odpowiada podobieństwu semantycznemu obiektów~\cite{hermans2017defense}. Kluczową rolę
w~tym procesie odgrywa uczenie metryk (\textit{ang. metric learning}) oparte na zoptymalizowanych
funkcjach straty~\cite{hermans2017defense}.

Wczesne architektury głębokie bazowały na sieciach siamskich (\textit{ang. Siamese networks})
składających się z~dwóch (lub więcej) bliźniaczych gałęzi o~współdzielonych wagach. Sieci te
optymalizowano przy użyciu funkcji straty kontrastowej (\textit{ang. contrastive loss}), która dąży
do minimalizacji odległości między cechami par o~tej samej tożsamości (pozytywnych) oraz
maksymalizacji odległości par o~różnych tożsamościach (negatywnych) powyżej zdefiniowanego
marginesu $\alpha$.

Powszechnie stosowaną alternatywą stała się funkcja straty trójkowej (\textit{ang. triplet loss}),
która operuje na zestawach złożonych z~próbki kotwiczącej ($x_a$), próbki pozytywnej ($x_p$) oraz
próbki negatywnej ($x_n$)~\cite{hermans2017defense}. Klasyczny triplet loss dąży do spełnienia
warunku:
\begin{equation}
  D(x_a, x_p) + m < D(x_a, x_n)
  \label{eq:triplet_loss}
\end{equation}
gdzie $D(\cdot, \cdot)$ oznacza euklidesową odległość w~przestrzeni cech, a~$m$ reprezentuje narzucony margines bezpieczeństwa~\cite{hermans2017defense}. Hermans i~in.~\cite{hermans2017defense} wykazali, że efektywność tej straty w~Re-ID ulega drastycznej poprawie przy zastosowaniu strategii próbkowania najtrudniejszych przykładów wewnątrz mini-batcha (\textit{ang. batch-hard triplet loss}):
\begin{equation}
  \mathcal{L}_{BH} = \sum_{i=1}^{P} \sum_{a=1}^{K}
  \left[ \max_{p \neq a} D(x_a, x_p) - \min_{j \neq i} D(x_a, x_n) + m \right]_+
  \label{eq:batch_hard_triplet}
\end{equation}
gdzie $P$ oznacza liczbę unikalnych tożsamości w~batchu, $K$ liczbę obrazów na każdą tożsamość, a~$[\cdot]_+$ reprezentuje funkcję $\max(0, \cdot)$~\cite{hermans2017defense}. Algorytm ten koncentruje proces optymalizacji wyłącznie na najmniej podobnych parach pozytywnych oraz najbardziej zbliżonych parach negatywnych, eliminując dominację łatwych przykładów~\cite{hermans2017defense}.

\subsection{Architektury splotowe (CNN) dedykowane do zadań Re-ID}
\label{subsec:cnn_reid}

Większość splotowych modeli reidentyfikacji osób wykorzystuje architekturę ResNet-50 jako sieć
bazową (\textit{ang. backbone}) ze względu na optymalny kompromis między dokładnością a~narzutem
obliczeniowym~\cite{ristani2016, zheng2016person}. Ze względu na specyfikę zadania, opracowano
dedykowane paradygmaty uczenia reprezentacji:
\begin{itemize}
  \item \textbf{IDE} (\textit{ang. ID-discriminative Embedding}): Model wprowadzony przez Zhenga i~in.~\cite{zheng2017person_wild}. Traktuje on proces treningowy jako problem klasyfikacji wieloklasowej, gdzie każda tożsamość w~zbiorze uczącym stanowi oddzielną klasę (klasyfikacja z~użyciem straty cross-entropy)~\cite{zheng2017person_wild}. Podczas wnioskowania odrzuca się warstwę klasyfikacyjną, a~wyjściowy wektor cech z~przedostatniej warstwy (GAP) służy bezpośrednio jako embedding reidentyfikacyjny~\cite{zheng2017person_wild}.
  \item \textbf{PCB} (\textit{ang. Part-based Convolutional Baseline}): Zaproponowany przez Suna i~in.~\cite{sun2018beyond} w~celu rozwiązania problemu utraty informacji przestrzennych w~globalnych deskryptorach. PCB rezygnuje z~globalnego uśredniania cech (GAP) na rzecz sztywnego podziału wyjściowej mapy splotowej na $T$ poziomych pasów (najczęściej $T=6$), z~których każdy reprezentuje określony fragment ciała (głowa, klatka piersiowa, biodra, uda, podudzia, stopy)~\cite{sun2018beyond}. Dla każdego pasa obliczany jest niezależny lokalny wektor cech, optymalizowany własną stratą klasyfikacyjną, co pozwala modelowi na ekstrakcję drobnych, lokalnych szczegółów wyglądu (\textit{ang. fine-grained features})~\cite{sun2018beyond}.
\end{itemize}

\subsection{Architektury oparte na Transformerach (SOTA)}
\label{subsec:transformers_reid}

Aktualny stan wiedzy (\textit{ang. State-of-the-Art}) w~dziedzinie Re-ID jest zdominowany przez
modele oparte na Vision Transformerach (ViT), które dzięki mechanizmowi samoatencji (\textit{ang.
  self-attention}) potrafią modelować globalne zależności dalekiego zasięgu bezpośrednio z~sekwencji
patchy obrazu, bez zniekształceń wprowadzanych przez operacje downsamplingu
splotowego~\cite{he2021transreid}.

Przełomową architekturą tej klasy jest model \textbf{TransReID} opracowany przez He
i~in.~\cite{he2021transreid}. TransReID wprowadza dwa rewolucyjne moduły zintegrowane ze strukturą
transformera:
\begin{itemize}
  \item \textbf{JPM} (\textit{ang. Jigsaw Patch Module}): Moduł ten dzieli zakodowane patche obrazu na kilka grup, dokonując na nich operacji przesunięcia (\textit{ang. shift}) i~losowego przemieszania (\textit{ang. patch shuffle}) przed podaniem ich do kolejnych warstw sieci~\cite{he2021transreid}. Wymusza to na transformerze naukę cech uniwersalnych i~zapobiega nadmiernemu skupieniu modelu na jednej, dominującej części sylwetki, co wydatnie zwiększa odporność na częściowe okluzje i~deformacje pozy targetu~\cite{he2021transreid}.
  \item \textbf{SIE} (\textit{ang. Side Information Embeddings}): Pozwala na bezpośrednie wstrzyknięcie nie-wizualnych informacji kontekstowych (takich jak identyfikator fizyczny kamery czy kąt widzenia obiektu) do wektorów wejściowych w~postaci uczalnych wektorów kodujących~\cite{he2021transreid}:
        \begin{equation}
          Z_0 = [x_{class}; E_{p_1}; \dots; E_{p_N}] + E_{pos} + E_{cam} + E_{view}
          \label{eq:sie}
        \end{equation}
        gdzie $E_{cam}$ i~$E_{view}$ eliminują błędy klasyfikacji wynikające ze specyfiki fizycznego usytuowania danej kamery lub kierunku poruszania się obiektu~\cite{he2021transreid}.
\end{itemize}

\subsection{Problem luki dziedzinowej i adaptacja bez nadzoru (UDA)}
\label{subsec:domain_gap_reid}

Fundamentalnym ograniczeniem wdrożeniowym systemów Re-ID jest występowanie luki dziedzinowej
(\textit{ang. domain gap})~\cite{ristani2016}. Model wykazujący precyzję Rank-1 rzędu 95\% na
zbiorze, na którym był trenowany (domena źródłowa), po uruchomieniu na obrazie z~innej,
niewidocznej wcześniej sieci kamer (domena docelowa) doświadcza drastycznego spadku dokładności
(często poniżej 40\% Rank-1)~\cite{ristani2016, ge2020mutual}. Zjawisko to wynika z~różnic
w~charakterystyce fotometrycznej sensorów, odmiennych temperatur barwowych oświetlenia naturalnego
i~sztucznego, tła oraz lokalnych dystrybucji odzieży~\cite{ristani2016, ge2020mutual}.

W~celu niwelowania tego problemu bez konieczności kosztownej i~czasochłonnej anotacji danych
docelowych rozwijane są metody nienadzorowanej adaptacji dziedzinowej (\textit{ang. Unsupervised
  Domain Adaptation} -- UDA)~\cite{ge2020mutual}. Wiodącym podejściem w~tej grupie jest algorytm
\textbf{MMT} (\textit{ang. Mutual Mean-Teacher}) zaproponowany przez Ge i~in.~\cite{ge2020mutual}.
MMT wykorzystuje architekturę nauczyciel-uczeń składającą się z~dwóch par identycznych sieci
szkolonych równolegle~\cite{ge2020mutual}. Modele nauczycieli są aktualizowane jako wykładnicza
średnia krocząca (\textit{ang. Exponential Moving Average} -- EMA) parametrów modeli
uczniów~\cite{ge2020mutual}. Poprzez generowanie miękkich pseudo-etykiet (\textit{ang. soft
  pseudo-labels}) w~domenie docelowej i~wzajemne nauczanie (\textit{ang. collaborative
  cross-teaching}), MMT skutecznie przeciwdziała akumulacji błędów klasyfikacji (szumu etykiet)
wynikającej z~błędnego klastrowania cech w~domenie nieznanej~\cite{ge2020mutual}.

\subsection{Specyfika reidentyfikacji pojazdów w odniesieniu do osób}
\label{subsec:vehicle_vs_person_reid}

Choć matematyczne podstawy reidentyfikacji pojazdów (\textit{ang. vehicle Re-ID}) i~osób są
tożsame, zadania te wykazują diametralne różnice o~charakterze metodycznym i~fizycznym:
\begin{itemize}
  \item Sylwetka człowieka jest wysoce odkształcalna (\textit{ang. deformable target}), co wymaga
        modelowania dynamicznych zmian pozy ciała, podczas gdy bryła pojazdu jest geometrycznie sztywna
        (\textit{ang. rigid target}).
  \item Samochody produkowane są seryjnie, co generuje problem skrajnie małych różnic międzyinstancyjnych
        (\textit{ang. small inter-instance variance}) -- dwa niezależne pojazdy tego samego modelu, koloru
        i~rocznika mogą być wizualnie identyczne, co wymusza na sieciach Re-ID koncentrację na unikalnych
        cechach personalizacyjnych, takich jak naklejki rejestracyjne, uszkodzenia karoserii czy przedmioty
        we wnętrzu kabiny.
  \item Wygląd pojazdu ulega drastycznym zmianom perspektywicznym pod wpływem rotacji (np. przód samochodu
        wykazuje zerowe podobieństwo wizualne do jego profilu), co wymaga stosowania podejść
        wielowidokowych (\textit{ang. multi-view Re-ID}).
\end{itemize}
Do ewaluacji algorytmów Re-ID pojazdów stosuje się dedykowane zbiory danych, z~których najważniejszymi są \textbf{VeRi-776} (zawierający ponad 50 000 zdjęć 776 unikalnych pojazdów zarejestrowanych przez 20 kamer) oraz \textbf{CityFlow} -- wielokamerowy benchmark o~skali miejskiej dedykowany do jednoczesnej detekcji, śledzenia i~reidentyfikacji w~trudnych warunkach ruchu drogowego~\cite{ristani2016}.

\subsection{Prywatność i alternatywne metody identyfikacji}
\label{subsec:privacy_reid}

Tradycyjne podejścia Re-ID bazujące na pełnozakresowych obrazach RGB sylwetek ludzkich budzą
poważne zastrzeżenia w~kontekście ochrony danych osobowych i~prywatności obywateli (np. w~świetle
unijnego rozporządzenia RODO)~\cite{ristani2016}. W~związku z~tym, nowoczesne badania koncentrują
się na metodach zapewniających prywatność z~poziomu samej architektury modelu (\textit{ang.
  privacy-by-design}):
\begin{itemize}
  \item \textbf{Silhouette Re-ID}: Identyfikacja prowadzona wyłącznie w~oparciu o~binarne maski sylwetek ruchu wycięte z~obrazu (np. przy użyciu modeli BGS). Metody te całkowicie eliminują teksturę ubrań i~twarz, opierając się wyłącznie na cechach antropometrycznych i~biomechanicznych chodu (\textit{ang. gait-based Re-ID})~\cite{chao2019gaitset}. Powszechnie stosowany w~tej dziedzinie model \textbf{GaitSet} traktuje sekwencję chodu jako nieuporządkowany zbiór sylwetek, co czyni go odpornym na ubytki klatek i~zmiany tempa poruszania się~\cite{chao2019gaitset}.
  \item \textbf{Skeleton-based Re-ID}: Podejście wykorzystujące algorytmy estymacji pozy (np. OpenPose) do wyznaczenia współrzędnych 3D kluczowych stawów anatomicznych człowieka~\cite{rao2023transg}. Modele takie jak \textbf{TranSG} (Rao i~in., 2023) reprezentują ruch człowieka za pomocą czasowo-przestrzennych grafów szkieletowych, wykorzystując transformery grafowe do uczenia unikalnych cech motorycznych przy zerowym dostępie do danych pikselowych twarzy czy odzieży, co gwarantuje pełną anonimizację wizerunku przy zachowaniu wysokiej stabilności śledzenia i~asocjacji cross-kamerowej~\cite{rao2023transg}.
\end{itemize}

\section{Rzutowanie na plan terenu (Bird’s-Eye View - BEV) i geometria sceny}
% To sekcja techniczna łącząca algebrę liniową z praktyką inżynierską — powinna stanowić pomost
% między kalibracją kamery (sekcja 1) a wielokamerowym śledzeniem (sekcja MTMC).

% Motywacja i cel rzutowania Problem przestrzeni obrazowej: współrzędne pikselowe z różnych kamer są
% nieporównywalne — ta sama lokalizacja fizyczna ma różne współrzędne ( u , v ) (u,v) w każdej
% kamerze; nie można wprost obliczać odległości ani prędkości

% Cel rzutowania na plan terenu: sprowadzenie wszystkich detektowanych obiektów do wspólnej
% metrycznej przestrzeni 2D (podłogi/terenu), gdzie odległości i pozycje są fizycznie
% interpretowalnie i porównywalne między kamerami

% Zastosowania w MTMC: fuzja trajektorii z wielu kamer, analiza spatio-temporalna przejść,
% wizualizacja sceny jako „mapa z lotu ptaka"

% Model geometryczny kamery Model otworkowy (pinhole camera model): rzutowanie punktu 3D współczynnik
% pochylenia (skew) — parametry unikalne dla każdej kamery, wyznaczane w procesie kalibracji

% Parametry zewnętrzne (extrinsic parameters): pozycja i orientacja kamery w przestrzeni świata —
% kluczowe dla rekonstrukcji przestrzennej i rzutowania na plan terenu

% Modele dystorsji: dystorsja radialna (barrel/pincushion) i tangencjalna — konieczna korekcja przed
% rzutowaniem; modele Browna-Conrady'ego

% Kalibracja kamery Metody kalibracji: szachownica (chessboard calibration, Zhang 2000) jako
% standardowe podejście; kalibracja za pomocą punktów referencyjnych w scenie (scene-based
% calibration) — praktyczna alternatywa dla kamer już zainstalowanych

% Kalibracja dla systemu CCTV: kamery stałe, więc kalibracja jednorazowa; problem zmiany parametrów w
% czasie (drgania, zmiany temperatury, wandalizm)

% Kalibracja wielokamerowa (multi-camera calibration): wyznaczenie wspólnego układu współrzędnych
% świata dla wszystkich kamer sieci; metody ze wspólnym polem widzenia i bez niego (non-overlapping
% FOV calibration)

% Narzędzia: OpenCV (calibrateCamera, stereoCalibrate), biblioteka Kalibr dla kamer wielokamerowych

% Homografia i rzutowanie perspektywiczne Definicja homografii: przekształcenie rzutowe między dwiema
% płaszczyznami opisane macierzą H H (3×3, 8 stopni swobody) — gdy płaszczyzna terenu jest płaska,
% rzutowanie obrazu na grunt jest izomorfizmem rzutowym

% Wyznaczanie macierzy homografii: minimum 4 niespółliniowe pary punktów korespondujących (image ↔
% world); algorytm DLT (Direct Linear Transform) z normalizacją; RANSAC do odrzucenia błędnych
% korespondencji (outliers)

% Transformacja odwrotna (BEV → obraz): cv2.warpPerspective — generowanie widoku z góry dla
% wizualizacji lub sieci neuronowych

% Ograniczenia homografii: zakłada płaski teren; błędy rosną dla obiektów wysokich (głowa osoby vs.
% stopy) — dyskusja o wyborze punktu referencyjnego (środek ciężkości, punkt kontaktu z podłożem)

% Wielowarstwowa homografia i rozszerzenia Problem wysokości obiektów: różne punkty na sylwetce leżą
% na różnych wysokościach nad gruntem → różne macierze homografii dla różnych warstw wysokości
% (multi-layer homography, np. Peng et al.)

% Podejście 3D (3D scene reconstruction): pełna rekonstrukcja 3D sceny z wielu kamer — bardziej
% dokładna, lecz kosztowna obliczeniowo; metody SfM (Structure from Motion), MVS (Multi-View Stereo)

% Probabilistyczna Mapa Zajętości (Probabilistic Occupancy Map): dla każdej komórki siatki na planie
% terenu oblicza się prawdopodobieństwo obecności obiektu jako iloczyn rzutowanych masek binarnych z
% wielu kamer — fundamentalne podejście Fleureta et al. (TPAMI 2008)

% Siatka BEV (BEV grid): dyskretyzacja planu terenu na komórki o rozmiarze np. 10×10 cm; kompromis
% między dokładnością a kosztem obliczeniowym

% BEV w nowoczesnych metodach głębokiego uczenia BEV jako wejście do sieci neuronowych: zamiast
% rzutowania geometrycznego — nauczalne moduły transformacji perspektywy; BEVNet, IPM (Inverse
% Perspective Mapping)

% Transformer-based BEV: BEVFormer (2022) — wielokamerowe BEV dla autonomicznej jazdy; architektura
% adaptowalna do stacjonarnych sieci kamer CCTV

% Zalety podejścia opartego na DL: odporność na niedoskonałości kalibracji, możliwość uczenia się na
% rzeczywistych danych, lepsza obsługa niespłaszczonego terenu

% Ograniczenia: potrzeba dużych zbiorów z anotacjami 3D lub BEV, trudna interpretowalność
% (explainability)

% Analiza spatio-temporalna w przestrzeni BEV Trajektoria w planie terenu: po rzutowaniu każdy
% tracklet z kamery staje się sekwencją punktów 2D w metrycznej przestrzeni; możliwa analiza
% prędkości, kierunku ruchu, stref wejścia/wyjścia

% Strefa wejścia/wyjścia (entry/exit zones): manualna lub automatyczna detekcja regionów, przez które
% obiekty wchodzą i wychodzą z pola widzenia kamery — kluczowe dla modelu przejść MTMC

% Predykcja pozycji między kamerami: jeśli znana jest mapa topologiczna sieci kamer, można szacować
% spodziewaną pozycję i czas pojawienia się obiektu w kolejnej kamerze na podstawie jego ostatniej
% pozycji w BEV

% Wizualizacja sceny: plan terenu jako interfejs operatora systemu — nakładanie trajektorii
% wszystkich obiektów na mapę lokalizacji (schemat budynku, zdjęcie lotnicze)

% Sugestia: tę sekcję warto umieścić po sekcji o kalibracji kamery (jako jej naturalne rozwinięcie),
% a przed sekcją MTMC — ponieważ rzutowanie na plan terenu jest technicznym fundamentem fuzji danych
% wielokamerowych. Możesz też zawrzeć tutaj własny schemat blokowy ilustrujący przepływ: obraz →
% detekcja bounding-box → punkt kontaktu z podłożem → homografia → pozycja BEV.

\section{Systemy wielokamerowe (MTMC) i fuzja danych}
% Definicja MTMC (Multi-Target Multi-Camera Tracking): problem utrzymania globalnych identyfikatorów
% obiektów w sieci kamer o pokrywających się lub niepokrywających polach widzenia (overlapping /
% non-overlapping FOV)

% Architektura systemu MTMC:

% Moduł detekcji (per-kamera)

% Lokalne śledzenie SCT (per-kamera, wynik: tracklets)

% Globalne skojarzenie między kamerami (inter-camera association)

% Zarządzanie globalnym ID

% Rzutowanie na plan terenu (ground plane projection): kalibracja kamery + homografia → wspólna
% metryczna przestrzeń 2D; probabilistyczna mapa zajętości (Probabilistic Occupancy Map, Fleuret et
% al., 2008); analiza spatio-temporalna przejść między kamerami

% Modele przejść między kamerami (camera transition model): rozkład czasu przejść (transition time
% distribution), topologia sieci kamer, strefy wejścia/wyjścia (entry/exit zones)

% Fuzja cech wizualnych i spatio-temporalnych: łączenie embeddingów Re-ID z informacją o czasie i
% przewidywanej pozycji — redukuje liczbę błędnych skojarzeń

% Podejścia grafowe i optymalizacyjne: graph-based multi-cut, sieć przepływów (flow network, Berclaz
% et al.), algorytmy klasteryzacji (hierarchical clustering, agglomerative clustering) trackletów

% Synchronizacja i kalibracja czasowa kamer: problem desynchronizacji strumieni wideo, timestamp
% alignment

% Benchmarki MTMC: DukeMTMC (wycofany z przyczyn etycznych), CityFlow (CVPR 2019), AICity Challenge —
% opis zbiorów i protokołów ewaluacyjnych

% Wyzwania otwarte: skalowalność na dziesiątki kamer, śledzenie w scenach z dużym tłumem,
% ograniczenia obliczeniowe dla systemów real-time, prywatność danych

\chapter{Przegląd istniejących rozwiązań}\label{chap:przegląd_rozwiązań}
\chapter{Implementacja systemu śledzenia i reidentyfikacji}\label{chap:implementacja}
\chapter{Integracja widoku wielokamerowego i rzutowanie na mapę}\label{chap:integracja_mapy}
\chapter{Badanie i analiza efektywności systemu}\label{chap:badanie_efektywności}
\chapter{Podsumowanie i wnioski}\label{chap:podsumowanie}

